\section{Motivasi}
Oke, ada $a^2+1$ terus kita mungkin tahu kalau $\tan^2 x + 1 = \sec ^2 x$. Terus di soalnya, $a,b,c$ bilangan real!! yaudah gasekun aja coba euy.

\section{Solusi}
Dikarenakan $a,b,c$ semuanya bilangan real dan fungsi $\tan$ hasilnya bisa mencakup semua bilangan real, maka masuk akal jika kita lakukan substitusi berikut.
\begin{align*}
a &= \tan x \\
b &= \tan y \\
c &= \tan z
\end{align*}
Berarti kita tinggal membuktikan ketaksamaan berikut
\begin{align*}
(ab+bc+ca-1)^2 &\leq (a^2+1)(b^2+1)(c^2+1)\\
(\tan x \tan y + \tan y \tan z + \tan z \tan x - 1)^2 &\leq (\tan^2 x + 1)(\tan^2 y + 1)(\tan^2 z + 1)\\
(\tan x \tan y + \tan y \tan z + \tan z \tan x - 1)^2 &\leq (\sec^2 x)(\sec^2 y)(\sec^2 z) \\
(\tan x \tan y + \tan y \tan z + \tan z \tan x - 1)^2& (\cos^2 x)(\cos^2 y)(\cos^2 z) \leq 1
\end{align*}
Di lain sisi kita punya
\begin{align*}
&(\tan x \tan y + \tan y \tan z)(\cos x \cos y \cos z) \\
&= \sin x \sin y \cos z + \sin y \sin z \cos x \\
&= \sin y (\sin(x+z))
\end{align*}
dan juga kita punya
\begin{align*}
&(\tan z \tan x - 1)(\cos x \cos y \cos z) \\
&= \sin x \sin z \cos y - \cos y \cos x \cos z \\
&= \cos y (-\cos(x+z))
\end{align*}
oleh karena itu ketaksamaan yang ingin dibuktikan setara dengan
\begin{align*}
&\left[ \sin y \sin(x+z) - \cos y \cos(x+z) \right]^2\\
&= \left[ -\cos(y+x+z) \right]^2 \\
&\leq 1
\end{align*}
yang jelas benar. Kesamaan dapat terjadi saat $a=b=c=0$.